\section{Contrastive loss objective derivation}\label{app:energy-refinement}
Given the model density (where $\beta_{\mathrm{mh}}$ is a constant):
\begin{equation}
p_\theta(x)=\frac{e^{-\beta_{\mathrm{mh}} V_\theta(x)}}{Z_\theta},
\qquad
Z_\theta=\sum_x e^{-\beta_{\mathrm{mh}} V_\theta(x)}.
\end{equation}
let us consider maximum likelihood estimation:
\begin{equation}
\theta^\star=\arg\max_\theta\;\mathbb{E}_{x\sim p_{\text{data}}}\big[\log p_\theta(x)\big]
\;\;\Longleftrightarrow\;\;
\theta^\star=\arg\min_\theta\;\mathcal{J}(\theta),
\end{equation}
with
\begin{equation}
\mathcal{J}(\theta):=-\mathbb{E}_{p_{\text{data}}}\big[\log p_\theta(x)\big]
=\mathbb{E}_{p_{\text{data}}}\big[\beta_{\mathrm{mh}} V_\theta(x)\big]+\log Z_\theta.
\end{equation}
Taking gradients gives
\begin{equation}
\nabla_\theta \mathcal{J}(\theta)
=\mathbb{E}_{p_{\text{data}}}\big[\beta_{\mathrm{mh}}\nabla_\theta V_\theta(x)\big]
+\nabla_\theta \log Z_\theta,
\end{equation}
where
\begin{equation}
\begin{aligned}
\nabla_\theta \log Z_\theta
&=\frac{1}{Z_\theta}\nabla_\theta Z_\theta
=\frac{1}{Z_\theta}\sum_x e^{-\beta_{\mathrm{mh}} V_\theta(x)}\big(-\beta_{\mathrm{mh}}\nabla_\theta V_\theta(x)\big)\\
&=-\mathbb{E}_{p_\theta}\big[\beta_{\mathrm{mh}}\nabla_\theta V_\theta(x)\big].
\end{aligned}
\end{equation}
Therefore, the gradient simplifies to
\begin{equation}
\boxed{
\nabla_\theta \mathcal{J}(\theta)
=\beta_{\mathrm{mh}}\left(\mathbb{E}_{p_{\text{data}}}\big[\nabla_\theta V_\theta(x)\big]
-\mathbb{E}_{p_\theta}\big[\nabla_\theta V_\theta(x)\big]\right)
}.
\end{equation}
This motivates the contrastive loss objective \(\mathcal{L}_{\mathrm{CL}}\) introduced in \eqref{eq:cl}, where we approximate the expectation over \(p_\theta\) using samples initialized either at uniform noise and data samples in 50\%/50\% proportions for general training, or exclusively at data samples (0\%/100\%) for the specialized training of data-initialized unconditional generation models. We then iteratively run a sequence of \((\text{greedy} + N_{\mathrm{CL}})\) proposals, explicitly detaching gradients flowing through the sampler transitions (i.e., Markov chain proposal and acceptance steps) to avoid backpropagation through the stochastic sampling process itself, ensuring stable estimation.
